\documentclass[11pt]{article}
% Preliminary draft (branch: augmented/correctness-and-paper).
% Target venue: AAAI AI for Social Impact (AISI) track, AAAI-27.
% For submission, switch the class to aaai2026.sty (two-column, 7 pages,
% double-blind) and move the author block into \nocopyright form.
\usepackage[margin=1in]{geometry}
\usepackage{amsmath,amssymb,amsthm}
\usepackage{booktabs}
\usepackage{hyperref}

\newtheorem{theorem}{Theorem}
\newtheorem{proposition}{Proposition}
\newtheorem{definition}{Definition}
\newtheorem{example}{Example}

\title{\textbf{Dynamic Welfare-Maximizing Pooled Testing\\with Count-Valued (Augmented) Outcomes}}
\author{Vladimir \and Hector \and Francisco Marmolejo-Coss\'io}
\date{Preliminary draft --- \today}

\begin{document}
\maketitle

\begin{abstract}
Pooled testing lets a public-health planner clear many people for work, school,
or travel with few tests, by testing samples together. Most welfare-oriented
work assumes a \emph{binary} pooled test: a pool is either all-clear or not. Yet
the laboratory signal is usually richer --- quantitative PCR returns a cycle
threshold that is, to first order, monotone in how much virus the pool carries.
We study the idealized limit of this richer signal: an \emph{augmented} pooled
test that reveals the exact \emph{count} of infected samples in the pool, and we
ask how this extra information changes \emph{dynamic} welfare-maximizing testing,
where each test is chosen after seeing the previous outcomes. We formalize the
problem, place its optimal value in a chain between single testing and the
full-information bound, and prove that exact posterior inference under count
outcomes is \#P-hard, while identifying the structured regimes (disjoint,
laminar, bounded-treewidth, few-test) that remain tractable. On the algorithmic
side we give an exact dynamic program, a myopic greedy policy, and a
reinforcement-learning baseline, and we correct two subtleties that otherwise
bias the evaluation: the outcome branches of a counting policy must be weighted
by the \emph{exact} conditional $P(r\mid\text{history})$, not by a
product-of-marginals approximation, and a Gibbs sampler over the exact-count
fiber must use moves that change the total infected count to be ergodic. On
synthetic instances the welfare advantage of the count test over the binary test
is small at low budget and grows with scale ($+0.6\%$ at $N{=}3$, $+4.0\%$ at
$N{=}5$, $+5.1\%$ at $N{=}7$), matching the intuition that count outcomes pay
off through better posteriors over a longer horizon.
\end{abstract}

\section{Overview}
A planner has a population of $N$ individuals. Each individual $i$ carries a
utility $u_i \ge 0$ --- the social value of confirming $i$ healthy so they may
return to in-person activity --- and a prior infection probability $p_i \in
[0,1]$, with $q_i = 1-p_i$. Infections are independent across individuals, a
modeling assumption appropriate to quarantine-style settings. The planner has a
budget of $B$ tests, each applied to a \emph{pool} of at most $G$ individuals,
where $G$ reflects a biological dilution limit. The planner earns $u_i$ once $i$
has been \emph{shown} healthy by a test. As in the static welfare-maximizing
framework of Finster et al. and its dynamic, binary extension by Lopez,
Marmolejo-Coss\'io, Tello Ayala and Parkes, the goal is to choose pools so as to
maximize the total utility of the individuals confirmed healthy by the end of the
protocol.

\paragraph{From binary to augmented tests.}
In classical pooled testing a single test on a subset $t \subseteq [N]$ returns a
binary answer: $-$ if nobody in $t$ is infected, and $+$ if at least one is. In
practice the test says more. A qPCR assay returns an integer in
$\{0,1,\dots,35,36\}$ (a cycle threshold) that is inversely related to the viral
load of the sample: a result of $36$ means no perceptible load (a negative), and
anything below means \emph{somebody} in the pool is infected. To begin
understanding what this extra signal buys, we study its idealized limit, where a
test on $t$ reveals the exact number of infected individuals it contains,
\[
  r(t,Z) \;=\; |t \cap Z| \;=\; \textstyle\sum_{i \in t} Z_i,
\]
with $Z \in \{0,1\}^N$ the (latent) infection profile, $Z_i \sim
\mathrm{Bernoulli}(p_i)$. We call such a test \emph{augmented}. The binary test is
the coarsening $\mathbf{1}[r > 0]$ of the augmented one, so the augmented setting
can only help; the question is by how much, and whether it can be exploited
efficiently. A balls-and-bins reading makes the idealization concrete: think of
each individual as a bag of red and white balls whose red fraction is the viral
load; a real assay reports a coarse range for the number of red balls, and the
augmented test is the limit in which that range collapses to the exact count.

\section{The Dynamic Augmented Pooled Testing Problem}
We follow the notation of the working note. A \emph{population instance} is
$J=(p_1,\dots,p_N,u_1,\dots,u_N)$. For a finite set $S$ and integer $G\ge 0$ let
$\mathcal{P}_G(S)=\{U\subseteq S : |U|\le G\}$.

\begin{definition}[Testing history]
A history of length $k$ is a list $h_k=\big((t_1,r_1),\dots,(t_k,r_k)\big)$ of
test--result pairs, $t_j\subseteq[N]$ and $r_j\in\{0,\dots,|t_j|\}$. The only
history of length $0$ is the empty one, $h_0=\emptyset$.
\end{definition}

\begin{definition}[DAPTS]
A \emph{dynamic augmented pooled testing strategy} for budget $B$ and pool limit
$G$ is a tuple $F=(F_1,\dots,F_B)$ with $F_k:\mathcal{H}_{k-1}\to
\mathcal{P}_G([N])$ mapping the observed history to the next pool.
\end{definition}

Running $F$ against a realized profile $Z$ produces a terminal history
$h_B(F,Z)$. The planner is paid for everyone placed in at least one all-clear
pool, so the realized utility is
\[
  u(F,Z)=\sum_{i\in[N]} u_i\,\mathbf{1}\!\left[\exists\,(t_k,r_k)\in h_B(F,Z):\,
  i\in t_k,\ r_k=0\right],
\]
and the value of $F$ is $u(F)=\mathbb{E}_{Z}\,u(F,Z)$, the expectation taken over
$Z\sim\prod_i \mathrm{Bernoulli}(p_i)$.

\paragraph{A hierarchy of optima.}
For a fixed $(J,B,G)$ it is useful to compare the best achievable value under
progressively more powerful regimes:
$\mathcal{U}^{\mathrm{single}}$ (test individuals one at a time, no pooling),
$\mathcal{U}^{s}_{\mathrm{NO}}$ and $\mathcal{U}^{s}_{\mathrm{O}}$ (optimal
\emph{static} non-overlapping and overlapping pooled plans, fixed in advance),
$\mathcal{U}^{D}$ (optimal \emph{dynamic binary} testing), $\mathcal{U}^{D}_{A}$
(optimal \emph{dynamic augmented} testing), and $\mathcal{U}^{\max}=\sum_i
u_i q_i$ (the value of knowing every status for free). By construction
\[
  \mathcal{U}^{\mathrm{single}} \le
  \mathcal{U}^{s}_{\mathrm{NO}} \le
  \mathcal{U}^{s}_{\mathrm{O}} \le
  \mathcal{U}^{D} \le
  \mathcal{U}^{D}_{A} \le
  \mathcal{U}^{\max}.
\]
The single new link is $\mathcal{U}^{D}\le\mathcal{U}^{D}_{A}$: the augmented
test dominates the binary one because its outcome refines the binary outcome.
The size of that gap --- the welfare value of counting --- is the quantity this
paper measures.

\begin{example}[Counting deduces what binary cannot]
Let $N=3$ with tests $t_1=\{0,1\}$ and $t_2=\{1,2\}$. Suppose the augmented
results are $r_1=1$ and $r_2=0$. From $r_2=0$ we learn individuals $1$ and $2$
are healthy; combined with $r_1=1$ we deduce individual $0$ is \emph{infected}
with certainty --- without ever testing $0$ alone. A binary planner seeing
$t_1{=}+,\,t_2{=}-$ learns the same about $1$ and $2$ but, in overlapping
designs with intermediate counts, the binary outcome generally loses the
cross-test arithmetic that pins down the remaining individuals.
\end{example}

\section{Posterior Inference under Augmented Pooled Testing}
Dynamic policies need posterior infection probabilities after each test. Writing
$A\in\{0,1\}^{k\times N}$ for the test--incidence matrix and $r$ for the observed
counts, the posterior over profiles is
\[
  \Pr[X=x \mid Ax=r] \;=\;
  \frac{\mathbf{1}[Ax=r]\,\prod_i p_i^{x_i}q_i^{1-x_i}}{Z},
  \qquad
  Z=\!\!\sum_{x:\,Ax=r}\prod_i p_i^{x_i}q_i^{1-x_i},
\]
and the planner wants the marginals $q_i'=\Pr[X_i=0\mid Ax=r]$. The normalizing
constant $Z$ is a weighted count of $0/1$ solutions to a system of cardinality
constraints, which is exactly where hardness enters.

\begin{proposition}\label{prop:sharpP}
Computing the exact posterior marginal $\Pr[X_i=0\mid Ax=r]$ under augmented
pooled testing is \#P-hard.
\end{proposition}
\begin{proof}[Proof sketch]
Reduce from \#Exact Cover. Associate a binary variable with each set in the set
system; for every ground element create a pooled test containing exactly the
variables whose sets contain that element, and fix its observed count to $1$.
Feasible $0/1$ assignments are then in bijection with exact covers, and with
uniform priors $p_i=\tfrac12$ the posterior is uniform over feasible
assignments, so $Z$ counts exact covers. Since counting exact covers is
\#P-complete, exact inference is \#P-hard.
\end{proof}

This is the central structural fact: the very deductions that make counting
powerful also make exact inference intractable in general. Several structured
regimes escape it.
\begin{itemize}
  \item \textbf{Disjoint pools.} If the tested subsets are pairwise disjoint the
  posterior factorizes and marginals are computed per pool.
  \item \textbf{Laminar / nested pools.} If every two pools are disjoint or
  nested, dynamic programming on the inclusion tree gives exact marginals.
  \item \textbf{Bounded treewidth.} The history induces a factor graph (variables
  = individuals, factors = cardinality-constrained tests); junction-tree
  inference is polynomial when its treewidth is bounded.
  \item \textbf{Few tests.} Grouping individuals by their membership pattern
  across the $k$ tests gives inference exponential in $k$ but polynomial in $N$.
\end{itemize}

\paragraph{Approximate inference and an ergodicity subtlety.}
Beyond these regimes we estimate marginals by sampling profiles consistent with
$Ax=r$. The natural Markov chain resamples one individual at a time, but on the
exact-count fiber a single-site move almost never preserves feasibility: any
agent inside a tight pool cannot flip alone. Swap moves (exchange an infected and
a healthy member of one pool) fix this within a pool but \emph{preserve the total
infected count}, so the chain is trapped in whichever count level it starts in.
On tests $\{0,1\}{=}1,\{1,2\}{=}1$ with prior $0.15$, the exact posterior is
$(0.15,0.85,0.15)$, yet a swap-only sampler returns $(0,1,0)$ or $(1,0,1)$
depending on the seed. We restore ergodicity in two steps. First we decompose
the active individuals into connected components (linked by shared tests); the
posterior factorizes across components, so each is solved independently and
\emph{exactly} by enumeration whenever it is small --- which, because the cap is
per component, covers every scale we encounter. For a single component too large
to enumerate we run a Metropolis sampler over \emph{alternating-path} moves:
kernel vectors of $A$ with entries in $\{-1,0,+1\}$ that balance every test yet
can change the total count (e.g.\ $(+1,-1,+1)$ on the example above), which is
precisely the move the local samplers lack.

\section{Algorithms}
\paragraph{Exact dynamic program.}
The optimal DAPTS is computed by a backward recursion over the state
$(k,\,\text{set of profiles still consistent with the history},\,\text{cleared
set})$. An augmented test branches into $|t|+1$ outcomes against the binary test's
two, and we reconstruct the optimal policy tree. Exact solution is feasible only
for small populations ($N\le 14$), which is why we also study scalable
heuristics.

\paragraph{Myopic greedy.}
At each step the greedy picks the pool maximizing the immediate expected cleared
utility,
\[
  \mathrm{Score}(t)=\Big(\textstyle\prod_{i\in t}(1-\tilde p_i)\Big)
  \sum_{i\in t} u_i,
\]
with $\tilde p_i$ the current posterior marginal, then updates beliefs and
repeats. A useful observation is that the myopic \emph{choice} is identical for
binary and augmented tests --- only $r=0$ yields immediate utility --- so the
augmented information helps only through better posteriors at later steps, never
through the first move. Two subtleties matter for correctness.
\begin{itemize}
  \item \textbf{Harvesting deduced-healthy individuals.} Utility is credited only
  when an individual is physically placed in an all-clear pool. An individual
  deduced healthy by counting ($\tilde p_i=0$) carries zero risk yet still owns
  $u_i$, which is collected only by testing them in a guaranteed-$r{=}0$ pool.
  Filtering such individuals out of future pools forfeits $u_i$ --- a genuine
  sub-optimality the dynamic program avoids by spending a test to harvest them.
  Keeping them eligible recovers the gap (on one instance, the counting greedy's
  value rises from $46.0$ to $57.6$, against the optimum $57.8$).
  \item \textbf{Exact branch weights in evaluation.} To compute a counting
  policy's expected value one sums over outcomes $r$. Weighting those branches by
  a Poisson-Binomial of the posterior \emph{marginals} assumes an independence
  that conditioning on the history has destroyed, and can mis-state a policy's
  own value by double digits. The branches must be weighted by the exact
  $P(r\mid\text{history})$ over consistent profiles; doing so makes each
  closed-form value equal the value obtained by simulating the policy.
\end{itemize}

\paragraph{Reinforcement learning.}
For populations beyond the exact wall we provide a reinforcement-learning
baseline: a bucketed Gymnasium environment (agents summarized by health/utility
histogram) in which a PPO agent assembles the first pool and the augmented myopic
greedy plays the remaining tests, with reward equal to the realized cleared
utility. The environment uses no conic solver and is fully seeded; training at
scale is left to a compute cluster.

\section{Experiments}
We draw $p_i\sim U(0,1)$ and $u_i\sim\mathrm{Uniform}\{1,2,3\}$ and report the
mean of each quantity in the hierarchy over $200$ instances per configuration
(corrected code; the chain held on every instance). Table~\ref{tab:hier} shows
the welfare advantage of the count test over the binary test, $\mathcal{U}^{D}_{A}
-\mathcal{U}^{D}$, growing from $+0.6\%$ at $N{=}3$ to $+4.0\%$ at $N{=}5$ and
$+5.1\%$ at $N{=}7$ (single-budget regimes show the smallest gap, as only one
outcome can be conditioned upon). The trend is consistent with the mechanism:
counting helps through posteriors over a longer horizon, so its value rises as
the budget and population grow.

\begin{table}[t]
\centering
\caption{Mean welfare (corrected code) over $200$ instances/config for $N\in
\{3,5\}$ and $40$ instances for the heavier $N{=}7$. The last column is the
augmented benefit $\mathcal{U}^{D}_{A}-\mathcal{U}^{D}$ as a percentage of
$\mathcal{U}^{D}$. The chain held on every instance.}
\label{tab:hier}
\small
\begin{tabular}{lllrrrrrr}
\toprule
$N$ & $B$ & $G$ & $\mathcal{U}^{\mathrm{single}}$ & $\mathcal{U}^{s}_{\mathrm{NO}}$ & $\mathcal{U}^{s}_{\mathrm{O}}$ & $\mathcal{U}^{D}$ & $\mathcal{U}^{D}_{A}$ & benefit \\
\midrule
3 & 2 & 3 & 2.572 & 2.643 & 2.659 & 2.686 & 2.704 & $+0.63\%$ \\
5 & 3 & 5 & 4.284 & 4.473 & 4.519 & 4.600 & 4.771 & $+3.97\%$ \\
7 & 3 & 7 & 4.908 & 5.249 & 5.335 & 5.454 & 5.717 & $+5.07\%$ \\
\bottomrule
\end{tabular}
\end{table}

\section{Conclusion and Outlook}
We have taken the step flagged as future work by the binary dynamic study ---
``richer testing outcomes beyond binary results'' --- and characterized dynamic
welfare-maximizing pooled testing under exact count outcomes. The picture has two
sides. Information: counting strictly dominates and its welfare value grows with
scale, but the inference that realizes it is \#P-hard in general and tractable
only on structured designs. Computation: a corrected greedy and an ergodic
sampler make the heuristics faithful, and a seeded RL environment is ready for
large-scale training. The open theoretical question is whether the count-augmented
per-step value is adaptive-submodular, which would give the greedy a
$(1-1/e)$-style guarantee; the open empirical question is how the advantage
behaves on real quantitative-PCR data, where the count is observed only through a
noisy cycle threshold. Both bear directly on whether augmented pooled testing can
improve real public-health screening.
\end{document}
